\chapter{What recent learning methods change, and what they inherit}
\label{r3:recent}

The preceding chapters distinguish the intervention, identification functional,
statistical experiment and numerical method. This distinction provides a way
to read the current literature without equating architectural novelty with a
new identification result. A transformer can estimate a conditional law, an
orthogonal learner can estimate a moment, and a simulator can integrate a
g-formula. They solve different parts of the same argument. For the bank,
the useful comparison is the scientific work each method removes, the
information it still requires, and the error it leaves in the financial target.

This chapter uses preserved primary versions available by 7 October 2026.
Recent preprints are research comparators. Their reported experiments are
not replications performed for this product, and a result for one observation
or assignment model is not transferred to a different bank experiment.

\section{Learning a trajectory distribution and learning its mean}
\label{r3:neural-g}

Write a recorded history as $H_t=(L_0,A_0,\ldots,L_t)$ and suppose the
conditions for Equation~\eqref{r3:sequential-g} hold. A generative route fits
the conditional laws $p_t(l_{t+1}\mid h_t,a_t)$ and simulates them under a
specified regime. G-Net implements this route with recurrent conditional
models and a within-visit factorization of the multivariate covariates
\citep[Sections 3--4 and Algorithm 1]{gnet2020}. If $L_{t+1}$ has $J$
ordered blocks, the factorization is
\[
 p_t(l_{t+1}\mid h_t,a_t)=
 \prod_{j=1}^J p_{tj}(l_{t+1,j}\mid h_t,a_t,l_{t+1,1:j-1}).
\]
This is a chain rule. Its within-visit order is not, without additional
assumptions, a learned causal order. An intervention on one of those
covariates therefore needs its own structural justification.

Under an adaptive regime $a_t=g_t(h_t)$, simulation must carry the generated
covariates into the next action. Reusing historical future actions mixes
the historical and proposed policies. Under a fixed future action sequence,
one may instead estimate the conditional mean of the final outcome directly.
These are different computational routes to an identified target; neither
changes sequential exchangeability or support.

Distributional simulation matters whenever a future mechanism is nonlinear.
For a project example, let a standardized balance obey
$L_1=a_0+\epsilon$, where $\epsilon$ is equally likely to be $-1$ or $1$,
and let a final funding-cost component be $Y=L_1^2$. Then
\[
 \E[Y^{a_0}]=a_0^2+1,\qquad
 \{\E[L_1^{a_0}]\}^2=a_0^2.
\]
For a rule that contacts a customer when $L_1<0$, at $a_0=0$ the contact
probability is $1/2$; feeding the mean zero into that rule gives zero.
Accurate one-step conditional means therefore do not suffice for general
forward simulation or threshold policies. Directly fitting the correctly
integrated terminal conditional mean is a possible alternative; it should
not be confused with plugging mean intermediate outcomes into a mechanism.

G-Net's continuous-variable implementation adds residuals sampled from a
holdout set to fitted conditional means \citep[Section 4, equation (5)]{gnet2020}.
That particular sampling scheme uses a residual law invariant to the
conditioning history. With history-dependent volatility, a correctly fitted
mean plus pooled residuals can still simulate the wrong tails and action
frequencies. Conditional scale/distribution models are a substantive
alternative. Simulation error, estimation error and counterfactual outcome
variation are distinct: generating more paths reduces Monte Carlo error,
but does not estimate the sampling uncertainty of the fitted mechanisms.

Causal Transformer targets conditional outcomes under specified future
treatment sequences, with consistency, sequential ignorability and overlap
\citep[Section 3 and Appendix A]{melnychuk2022}. It uses attention over
observed histories and autoregressively feeds predictions. The example
above is a qualification test for any interpretation of such predictions
as a generative simulator; it is not a claim to have refuted that paper's
reported benchmark results. A causal attention mask prevents certain
information leaks. It cannot supply an omitted confounder or identify an
unsupported treatment sequence.

\section{Balancing representations needs a sufficiency argument}
\label{r3:balance-limit}

Counterfactual Recurrent Networks train a history representation
$R_t=f(H_t)$ using outcome prediction and an adversarial treatment
classifier. Theorem 1 of \citet{bica2020} characterizes equality of the
treatment-group representation distributions at the ideal optimum of
the treatment-adversarial game. That theorem alone does not establish
\[
 Y^g\perp A_t\mid R_t.
\]
Sequential ignorability given $H_t$ need not survive its compression.
Equality $P(R_t\mid A_t=a)=P(R_t\mid A_t=b)$ is also different from the
balancing-score property $A_t\perp H_t\mid R_t$, which, together with the
original ignorability condition, can justify adjustment.

Consider the one-period special case $H\sim\operatorname{Bernoulli}(1/2)$,
$P(A=1\mid H)=.2+.6H$, and $Y=A+H$. All relevant confounding is recorded
and the treatment effect is one. The constant representation $R=0$ has
identical distributions in both treatment groups, yet
\[
 \E[Y\mid A=1,R]-\E[Y\mid A=0,R]
 =1+\{.8-.2\}=1.6.
\]
An outcome-prediction term may discourage this particular representation;
the counterexample does not describe every minimizer of the combined
training objective. It establishes the narrower logical point that
distributional balancing is insufficient for identification.

A justified representation can be invertible, can retain a valid balancing
score, or can satisfy target-specific outcome/transition sufficiency. For
dynamic policies it must also preserve the information needed to choose
actions and integrate future outcomes. A representation trained for
one-period response prediction need not be sufficient for long-horizon
policy value. Adversarial accuracy, reconstruction error and effect
calibration probe different failures and should be reported separately.

\section{Orthogonal dynamic learning is a moment construction}
\label{r3:recent-dynamic-dml}

The recursive blip argument in Section~\ref{r3:msm-snm} leads to a different
use of machine learning. In the linear dynamic setting of
\citet[Section 2, Theorem 2 and Algorithm 1]{dynamicdml2021}, let $A_t$
be the treatment vector, $X_t$ the sufficient recorded state, $Y$ the
terminal outcome, and $\psi_t$ the effect of treatment at $t$ relative
to the specified future reference treatment. Define nuisance means
\[
 q_t(X_t)=\E[Y\mid X_t],\qquad
 p_{jt}(X_t)=\E[A_j\mid X_t],\quad j\geq t.
\]
Residualize $\widetilde Y_t=Y-q_t(X_t)$ and
$\widetilde A_{jt}=A_j-p_{jt}(X_t)$. Backward estimation solves
\begin{equation}
 \E\left[\widetilde A_{tt}
 \left\{\widetilde Y_t-\sum_{j=t}^T\psi_j^\top
                  \widetilde A_{jt}\right\}\right]=0.
 \label{r3:dynamic-orthogonal}
\end{equation}
The terminal equation identifies $\psi_T$ if the residual treatment
covariance is invertible. Subtracting that blip exposes the previous
equation. This is the same causal peeling operation as g-estimation,
with nuisance regression organized to make its moment orthogonal.

For example, perturbing $q_t$ by a function of $X_t$ has zero first-order
effect because $\E[\widetilde A_{tt}\mid X_t]=0$. Perturbing a
$p_{jt}$ has the same cancellation; perturbing the leading treatment
residual also multiplies a mean-zero residual conditional on $X_t$ at
the truth. Finite nuisance errors leave products, while an error in a
later $\widehat\psi_j$ propagates backward through cross-treatment moments.
Consequently orthogonality to nuisance regressions does not remove
horizon-dependent conditioning or the need for joint covariance inference.

The initial theorem uses independent trajectories, fixed low treatment
dimension and horizon, a specified linear structural model and nonsingular
residual treatment variation. The source separately treats more general
blips, heterogeneity and a dependent-series setting. These are distinct
theorem regimes. A package that accepts a panel identifier does not thereby
inherit all of them, and nonlinear nuisance learners do not make an
incorrectly linear treatment-response restriction disappear.

Even ordinary panel residualization needs care.
\citet[Sections 2--4]{paneldml2024} compares early/late demeaning and
correlated-random-effect adaptations. If the true nuisance is
$f(X_{it})=X_{it}^2$, then its within transformation is
\[
 X_{it}^2-\overline{X_i^2}
 =2\bar X_i\widetilde X_{it}
   +\widetilde X_{it}^2-\overline{\widetilde X_i^2}.
\]
Two units with paths $(0,2)$ and $(10,12)$ have identical demeaned
covariates $(-1,1)$, but demeaned quadratic responses $(-2,2)$ and
$(-22,22)$. Learning only from the demeaned covariate loses needed
information. Adding history summaries can help under an explicit
correlated-effect model; it is not a general proof of removal of latent
heterogeneity. The inspected panel comparison emphasizes point-estimator
simulation, so general dependent-data confidence intervals are not inferred
from that comparison.

\section{Amortized causal inference inherits an identification model}
\label{r3:causal-fm}

Let $M$ be a structural model drawn from a prior $\Pi$, $P_M$ its
observational law, and $\theta_M$ the requested causal functional.
A prior-data fitted network can minimize the population loss
\[
 \E_{M\sim\Pi,D_n\sim P_M^n}
       [-\log q_\eta(\theta_M\mid D_n)].
\]
With an unrestricted conditional density family and exact optimization,
the optimum is the prior's conditional target distribution. A finite
network trained on finitely many simulations only approximates that
optimum. If training labels are future outcomes rather than causal
functionals, the optimal distribution is different.

CausalFM constructs priors on compatible observational/interventional
distribution pairs, using practitioner-supplied back-door, front-door or
IV identification settings \citep[Definitions 3.1 and 4.2]{causalpfn2026}.
Its IV example uses an additive response and completeness conditions;
it is not an amortized solution for every valid binary-IV LATE problem.
The paper's Theorem 4.3 gives a prior-consistency necessity result for
a linear query, with its stated noncancellation and observational-class
conditions. Identification is not a sufficient guarantee that an arbitrary
prior or trained network is calibrated for every new dataset.

The finite case exposes the central restriction without an asymptotic
argument. The two models in Section~\ref{r3:nonidentification} give exactly
the same strictly positive distribution of $(A,Y)$, but have ATEs zero
and $.36$. If the second model has prior mass $w$, then, for every
observational sample,
\[
 P(M=2\mid D_n)=w,\qquad
 \E[\theta_M\mid D_n]=.36w,\qquad
 \operatorname{Var}(\theta_M\mid D_n)=.36^2w(1-w).
\]
The likelihood cannot revise the relative weights of equivalent models.
More training examples can approximate this prior conclusion more
accurately; more observations from the same experiment cannot identify
which causal model generated them. Additional causal information changes
the experiment or the allowed model class.

The same example bounds a learned classifier of those causal explanations.
Under equal prior probability, any classifier based only on the observed
sample has average classification error at least $1/2$, since its input
distribution is identical in the two cases. Identifiability itself is a
property of a target and an assumed model class, not a directly observed
label of one distribution. An output score can be meaningful under its
training prior; it cannot be a distribution-free certificate that a
bank's assumptions hold.

The temporal PFN preprint \citet[Sections 3.2--3.6]{panelpfn2026} proposes
time encodings, a causal-judgment head, mixed structural priors, selected
contexts and one-step posterior correction. These are useful research
questions, but a numerical identifiability score cannot defeat the
equivalence example. Nor does writing
$\widehat\psi_{\rm PFN}+\mathbb P_n\widehat\phi$
specify valid inference: one must identify the target, derive its
influence function for the actual longitudinal observation model, control
the remainder, account for context selection and dependence, and justify
the covariance estimate. The paper's displayed generic correction is
not accepted here as a theorem covering arbitrary panel CATEs.

\section{Proximal foundation models still face an inverse problem}
\label{r3:proximal-fm}

ProximalFM targets $\tau(x)=\E[Y(1)-Y(0)\mid X=x]$ by averaging over
the proxies rather than conditioning its target on them
\citep[Sections 2--3]{proximalfm2026}. Its causal route remains the bridge
equation in Section~\ref{r3:proxy-identification}, together with latent
exchangeability, proxy exclusions, positivity, completeness and bridge
existence. The network amortizes inference across synthetic structural
models; it does not make the integral operator injective.

The distinction is acknowledged in the primary text. Proxy exclusions
follow from its chosen structural inputs and independent noises, and
propensity clipping enforces positivity in the simulator. General nonlinear
completeness is not guaranteed. Proxy width and cardinality constraints
are used as heuristics, with relevance diagnostics
\citep[Section 3.1 and Appendix B.3.4]{proximalfm2026}.

For a finite illustration, let $U=(U_1,U_2)$ contain two independent
binary coordinates. Construct five proxy bits as conditionally independent
noisy measurements of $U_1$, with error probability $.2$. The proxy
vector has dimension five and 32 possible values, exceeding the latent
dimension two and four possible latent values. Nevertheless, its
$32\times4$ conditional-probability matrix has duplicate columns for
$(U_1,0)$ and $(U_1,1)$ and rank two. Width and cardinality do not
recover the missing coordinate. To make the conditional argument explicit,
let treatment be assigned from $U$ with strictly positive probabilities
using noise independent of the proxy noises. Writing the proxy vector as
$Z$, this gives $Z\perp (U_2,A)\mid U_1$. For each treatment stratum $a$,
the nonzero function
\[
 h_a(U)=U_2-\E[U_2\mid U_1,A=a]
\]
satisfies $\E[h_a(U)\mid Z,A=a]=0$ by iterated expectation. Centering only
on $A$ would in general be insufficient: conditioning on treatment can
associate $U_1$ and $U_2$. Thus incompleteness persists after conditioning
on treatment under the specified assignment law.
The relevant bridge and target must be examined as in
Equation~\eqref{r3:null-functional}; even when a particular target
survives a null space, that requires an argument.

The preprint's Appendix A carefully separates the oracle from its trained
network. Let $R(M)=P_M$ and fix a query $x$. Under measurability and
integrability, the oracle posterior mean converges, for $\Pi$-almost every
model, to
\[
 g_x(R(M))=\E_\Pi[\tau_M(x)\mid R(M)].
\]
Its Theorem A.4 makes oracle consistency equivalent to the target being
a measurable function of $R(M)$ almost surely under the prior. This is
prior-almost-sure, pointwise-query consistency. It is not uniform coverage
over bank models outside the prior, and the completeness heuristic does
not establish the theorem's identification premise for every generated SCM.

Training labels introduce another distinction. If
$\widehat\tau_K(x)=K^{-1}\sum_k\{Y_k(1)-Y_k(0)\}$ uses independent
downstream simulations conditional on the SCM and query, then
\[
 \E[\widehat\tau_K(x)\mid D_n]
       =\E[\tau_M(x)\mid D_n].
\]
Thus population squared-error training targets the oracle posterior mean
even at finite $K$. The full label distribution is convolved with simulation
noise. Recovering a posterior distribution needs an identified noise
model or vanishing label noise. Appendix A.4 explicitly does not prove
that the finite-$K$ Gaussian noise correction yields exact recovery.
Individual-effect variation, oracle posterior uncertainty, approximation
error and frequentist coverage are four different quantities.

\section{History-dependent logging can make informative paths rare}
\label{r3:ope-hardness}

A hidden-state model with few states can have a difficult value-estimation
problem even when actions are randomized conditional on recorded history.
\citet[Sections 2--5]{opehard2026} constructs two finite-horizon models
that differ only in the hidden state reached by the first continue
action. Later continue actions preserve the state; either of two reset
actions fixes it to zero or one. Only the terminal observation reveals
anything, with reward-one probabilities $1/4$ and $3/4$ in the two states.
The target always continues, so its values are $1/4$ and $3/4$.

The known logger chooses the first action uniformly. Before any reset it
continues with probability $2/3$, while reset probabilities are $1/6$
each. After a reset its history-dependent action probabilities change,
but every action remains at least $1/6$. A reset erases the only difference
between the models. A trajectory is informative precisely when it never
resets, with probability
\[
p_H=\frac13\left(\frac23\right)^{H-2},\qquad H\ge2.
\]
Reduce each trajectory to an erasure or its terminal bit. The two laws,
in the order $(\bot,0,1)$, are
\[
 Q_0=(1-p_H,3p_H/4,p_H/4),\qquad
 Q_1=(1-p_H,p_H/4,3p_H/4).
\]
The rest of the trajectory conditional on that symbol is distributed
identically in both models. Direct calculation gives
$D_{\rm KL}(Q_0\Vert Q_1)=p_H\log(3)/2$. If an estimator is within $1/8$
of the value with probability $1-\delta$ in both models, for
$0<\delta<1/2$, thresholding at
$1/2$ yields a test with both errors at most $\delta$. Data processing
of relative entropy through that test therefore gives
\begin{equation}
 n\geq\frac{2\,\operatorname{kl}(1-\delta,\delta)}
                  {p_H\log3}.
 \label{r3:rare-path-bound}
\end{equation}
No faster optimizer changes this missing information. For the reduced
experiment, the exact optimal test chooses model 1 when terminal ones
outnumber zeros among nonerased episodes, randomizing at a tie.

The striking part of the source theorem is that its particular
behavior-marginal state/belief/revealing coverage constants remain
bounded as $H$ grows. The future distribution in that definition averages
over logger memory conditional on physical state; it need not equal
the continuation distribution after a fixed history. Augmenting state
by logger memory repairs the Markov description but reveals a rare
never-reset state and a growing coverage cost. These are not hidden
confounding and not a claim of hardness for every POMDP. They show why
the conditioning set in an OPE theorem must be preserved by an adapter.

In banking terms, many customers may visit every coarse delinquency band,
yet very few follow the uninterrupted action path of a proposed policy.
The synthetic construction does not quantify that frequency in a bank.
It motivates measuring target-relevant trajectory support and considering
a shorter target horizon, additional information, or new randomized data
when the historical experiment is uninformative.

\section{Dynamic bounds require a compatible joint causal model}
\label{r3:dynamic-bounds}

Partial identification can replace a point forecast, but local uncertainty
sets must correspond to a common structural model.
\citet[Sections 2--4]{sharpdynamic2026} studies a finite observed-state
model with stage-specific hidden confounding between $A_t$ and $S_{t+1}$.
Its hidden disturbances do not persist as an unrecorded customer
characteristic across stages. At a confounded supported cell $(t,s,a)$,
write
\[
 m_{s'}=P(S_{t+1}=s',A_t=a\mid S_t=s),\qquad
 r=P(A_t\ne a\mid S_t=s).
\]
The default feasible next-state law is
\[
 {\cal Q}_t(s,a)=
 \{q\in\Delta({\cal S}_{t+1}(s,a)):m_{s'}\leq q_{s'}
                                  \leq m_{s'}+r\}.
\]
Equivalently, $q=m+r v$ with $v$ a probability vector on the allowed
successors (with the evident singleton when $r=0$). The observed
treated mass is fixed; the untreated counterfactual mass is unobserved.
Without support the set is the whole structural successor simplex;
with exchangeability it reduces to the observed transition.

For a fixed regime $d$ and terminal reward $y(s)$, rectangularity permits
the lower recursion
\[
 \underline V_T(s)=y(s),\qquad
 \underline V_t(s)=\min_{q\in{\cal Q}_t(s,d_t(s))}
                           \sum_{s'}q(s')\underline V_{t+1}(s'),
\]
and similarly a maximum for the upper bound. The proof proceeds backward:
one can choose an attaining local kernel at every state while keeping
all already chosen continuation kernels. Compactness gives attainment,
and the ability to paste those choices gives equality with the global
optimization. Without that ability, the recursion optimizes an enlarged
model class.

For example, if $m=(.2,.4)$, $r=.4$, and the next-state values are
$(0,1)$, the sharp local mean interval is $[.4,.8]$.
The source's Theorem 3.8 proves global causal sharpness for its
\emph{unmodified default} local sets and maintained causal model.
Narrowed sensitivity sets, estimated confidence regions or arbitrary
POMDP ambiguity sets do not automatically inherit that conclusion.

A project example makes the coupling problem concrete. Suppose a persistent
unobserved binary characteristic $U$ produces two cash-flow components
$C_0=U$ and $C_1=1-U$, with unknown distribution of $U$. Each marginal
mean may range over $[0,1]$, but $\E[C_0+C_1]=1$ under every allowed
model. Replacing the coupled model by the Cartesian product of the two
marginal intervals gives $[0,2]$. Such bounds are conservative but
not sharp for the specified model. Including sufficient state and
preserving the structural successor restriction can restore the relation;
independently optimizing each marginal cannot.

The 2026 paper takes the observational law as given. Finite-sample
uncertainty in $m$, $r$ and the support is another problem. Estimating them
and treating the resulting polytope as known understates uncertainty.
Joint confidence regions, moment-inequality inference or a justified
bootstrap for the bound endpoints need their own regularity analysis.

\section{Irregular discovery and deconfounding must match the observation law}
\label{r3:recent-discovery}

The ITSY preprint addresses missing entries on an underlying regular grid
in a linear first-order SVAR,
\[
 X_t=AX_t+BX_{t-1}+e_t,\qquad
 X_t=(I-A)^{-1}BX_{t-1}+(I-A)^{-1}e_t,
\]
with a contemporaneous DAG, causal sufficiency, structural stationarity
and MCAR missingness \citep[Sections 3--4]{itsy2026}. Reformulation
transforms the noise covariance:
\[
 \Omega_{\rm red}=(I-A)^{-1}\Omega_e(I-A)^{-\top}.
\]
Thus minimizing an unweighted error in the reduced equation can change
the structural objective. The source couples imputation and graph
learning and corrects the residual geometry. This is a meaningful
algorithmic contribution inside its model; it does not turn endogenous
customer visits into MCAR grid omissions. Continuous-time monitoring
requires the event-likelihood analysis of Chapter~\ref{r3:continuous}.

LUCID routes a discovery procedure between sparse and pervasive
confounding using the residual spectrum of a fitted VAR(1)
\citep[Section 4]{lucid2026}. Its branches use post-double-selection
filtering or spectral trimming and contemporaneous recovery. The
inspected text explicitly calls its lag-zero orientation step a
heuristic; a router-separation proof is referred to an extended version.
A low-rank residual component can reflect common causes, a misspecified
dynamic model, or common measurement effects. Removing a spectral
component changes the information used for discovery; it does not
generally prove that the remaining graph is causally sufficient.

These proposals should be tested against the distinctions already
derived: hidden common causes versus mediators, full and summary graphs,
lag errors versus direction errors, MCAR versus informative observation,
and observational equivalence versus structural identification. Synthetic
edge-recovery scores establish performance under their generating regimes.
For a bank decision the final question remains whether the graph class,
including unresolved orientations and latent structure, identifies the
requested intervention or only a set of possible effects.

\section{Comparison at the level of scientific operations}
\label{r3:recent-contract}

The promising reusable operations are conditional-law learning,
orthogonal moment solving, prior-conditioned amortization, constrained
graph search and robust value recursion. Each returns a different
mathematical object. A graph distribution is not a causal-effect confidence
interval; a posterior under an SCM prior is not a test of that prior's
identifying restrictions; a Bellman bound over a rectangular uncertainty
class is not necessarily a sharp causal bound.

The appropriate comparison holds the target, information set and sampling
unit fixed. For neural g-computation, compare simulated intervention
functionals and tail-dependent decisions as well as prediction error.
For orthogonal learners, inspect moment residuals, nuisance products,
split validity and weak conditioning. For foundation models, vary
structural priors, context selection, proxy strength and departures from
the training experiment. For discovery, retain equivalence information
and propagate it to effect sets. For policy evaluation, inspect
target-relevant paths and logger memory. This leads directly to a
software contract: an estimator can be interchangeable only within a
common, justified mathematical operation.
